\section*{Chapter \ref{ch:m2:credit-default-swaps} --- Credit Default Swaps}

\begin{solution}{exo:m2:credit-default-swaps:1}
USD~25\,000 a quarter ($10 \text{ million} \times 1\% / 4$); after the default, $10
\text{ million} \times (100 - 40)\% = \text{USD}~6$ million, and the premiums stop.
\end{solution}

\begin{solution}{exo:m2:credit-default-swaps:2}
$\lambda \approx 0.03 / 0.6 = 5\%$ a year; $Q(5) = e^{-0.25} = 77.9\%$. It is a
risk-neutral probability, not a forecast.
\end{solution}

\begin{solution}{exo:m2:credit-default-swaps:3}
Whether a \omterm{def:m2:credit-default-swaps:event}{credit event} has occurred, and when; whether to hold an auction; which
obligations are deliverable; whether the entity has a successor, and which. It also
answers questions of interpretation of the contract.
\end{solution}

\begin{solution}{exo:m2:credit-default-swaps:4}
The upfront is $-2.21\%$: the \omterm{def:m2:credit-default-swaps:cds}{protection buyer} receives about USD~221\,000, because it
will pay a coupon of 100 basis points for protection worth 50.
\end{solution}

\begin{solution}{exo:m2:credit-default-swaps:5}
Netting needs identical contracts: two trades on the same name and maturity but with
different coupons leave a residual stream of premiums that cannot be cancelled. With a
\omterm{def:m2:credit-default-swaps:std}{standard coupon} the difference is paid once, upfront, and the running legs are
identical, so opposite trades cancel and a central counterparty can hold one net
position per name and maturity.
\end{solution}

\begin{solution}{exo:m2:credit-default-swaps:6}
It delivers the bond at 30 and receives par: USD~10 million for bonds worth USD~3
million, a payout of USD~7 million. The auction's final price should reflect the
cheapest deliverable, about 30, since sellers of bonds in the auction deliver it.
\end{solution}

\begin{solution}{exo:m2:credit-default-swaps:7}
1.6667\%, 5.0001\% and 13.335\%, against 1.6667\%, 5.0000\% and 13.333\% from the
triangle. The small gap, growing with the spread, comes from paying premiums and
defaults quarterly rather than continuously; counting half a period of accrued
premium makes the discrete legs nearly match the continuous ones.
\end{solution}

\begin{solution}{exo:m2:credit-default-swaps:8}
The spread is about the \omterm{def:m2:credit-default-swaps:hazard}{hazard rate} times the loss given default, $\lambda(1-R)$: at
300 basis points and 40\% recovery the \omterm{def:m2:credit-default-swaps:hazard}{hazard rate} is about 5\% a year, not 3\%. And it
is a risk-neutral rate, which includes a premium for bearing default risk and
illiquidity; actual default rates are usually lower.
\end{solution}

\begin{solution}{pb:m2:credit-default-swaps:1}
\textbf{1.} The highest bid, 10, and the lowest offer, 10: they cross and are removed.
\textbf{2.} Seven of the thirteen remaining: bids 9.5, 9.5, 9.25, 9.25, 9, 9, 8.875 and
offers 10, 10, 10, 10.25, 10.75, 10.875, 11 average 9.804, rounded to 9.75.
\textbf{3.} Sells USD~5\,692 million, buys USD~772 million: USD~4\,920 million to sell.
\textbf{4.} The dealer that bid 10, HSBC, above the midpoint while the open interest was to
sell: $5\text{ million} \times (10 - 9.75)\% = \text{USD}~12\,500$.
\textbf{5.} \omterm{def:m2:credit-default-swaps:cds}{Protection buyers} that held bonds asked to sell them at the final price, to
end with par, and they outweighed \omterm{def:m2:credit-default-swaps:cds}{protection sellers} asking to buy.
\textbf{6.} Limit bids to buy bonds, at no more than $9.75 + 1 = 10.75$.
\textbf{7.} Cumulative sizes 50, 200, 500, 1\,400 and 2\,800 million down to 8.75, then
5\,400 million at 8.625, which meets 4\,920: the final price is 8.625.
\textbf{8.} 10.75, the cap.
\textbf{9.} 3\,800 million at 8.625 is not enough; with the 1\,800 million at 8.5 the
book reaches 5\,600: the final price is 8.5.
\textbf{10.} 1.125 below the midpoint: the open interest to sell was large, and the bids
needed to absorb it were lower.
\textbf{11.} $20 \text{ million} \times 91.375\% = \text{USD}~18.275$ million; its bonds
are worth USD~1.725 million at 8.625.
\textbf{12.} USD~18.275 million from the CDS and USD~1.725 million for the bonds: USD~20
million, par, and no bonds left.
\textbf{13.} USD~45.6875 million.
\textbf{14.} On its net USD~40 million sold, USD~36.55 million.
\textbf{15.} USD~27.4125 million on the CDS and USD~2.5875 million for the bonds: USD~30
million in all, for USD~30 million face of bonds, as if it had been delivered them.
\textbf{16.} Most of the protection was held by dealers that had both bought and sold
it; in the auction each paid or received only its net position.
\textbf{17.} The open interest to sell was USD~4.92 billion; absorbing that supply took
bids below the bond market's last prices.
\textbf{18.} Holders of protection without bonds would have had to buy them, pushing the
price up and their payout down, and contracts on the same name would have settled at
different prices.
\textbf{19.} Named result: \emph{the auction} sets the final price at 8.625; A receives
USD~18.275 million and keeps bonds worth 1.725 (or USD~20 million with no bonds), B
receives USD~45.6875 million, C pays USD~36.55 million, and D pays USD~30 million in all
for USD~30 million face of bonds.
\textbf{20.} It gives every contract on a name one recovery price, set by trading the
net demand for its bonds, so settlement no longer depends on finding bonds.
\end{solution}

\begin{solution}{iq:m2:credit-default-swaps:1}
A contract in which the \omterm{def:m2:credit-default-swaps:cds}{protection buyer} pays a running premium on a notional to the
seller, quarterly, until maturity or a \omterm{def:m2:credit-default-swaps:event}{credit event} of the \omterm{def:m2:credit-default-swaps:cds}{reference entity}; after a
\omterm{def:m2:credit-default-swaps:event}{credit event} the seller pays the notional times one minus the recovery price set by
the auction. The buyer need not own the debt.
\iqlookfor{the two legs, the trigger, settlement.}
\end{solution}

\begin{solution}{iq:m2:credit-default-swaps:2}
Since 2009 contracts carry a fixed coupon (100 or 500 basis points for North American
names) and the difference from the market spread is paid upfront. Identical running
legs make trades on the same name and maturity fungible: they can be netted,
compressed and cleared.
\iqlookfor{fungibility, netting, clearing.}
\end{solution}

\begin{solution}{iq:m2:credit-default-swaps:3}
With a flat hazard $\lambda$, recovery $R$ and continuous premiums, the premium leg is
$s\int e^{-(r+\lambda)t}dt$ and the protection leg $(1-R)\lambda\int
e^{-(r+\lambda)t}dt$; equating gives $s = \lambda(1-R)$. Discrete payments change it
slightly; the hazard is risk-neutral.
\iqlookfor{the two legs and the cancellation.}
\end{solution}

\begin{solution}{iq:m2:credit-default-swaps:4}
First stage: dealers submit bids and offers for a fixed size, and physical settlement
requests; crossing pairs are removed and the better half averaged into the inside
market midpoint; the net requests are the open interest. Second stage: limit orders
on the other side fill the open interest, best first, and the last order filled sets
the final price, within one point of the midpoint. It is needed because protection can
exceed deliverable bonds and bonds trade thinly after default; it gives one price for
all contracts.
\iqlookfor{both stages, the cap, the reason.}
\end{solution}

\begin{solution}{iq:m2:credit-default-swaps:5}
The basis is the CDS spread minus the bond's spread over the funding rate (for a
floating-rate bond, its spread; for a fixed bond, an asset-swap or Z-spread). If
negative, buy the bond, fund it in repo and buy protection: you earn the difference and
are hedged against default. The risks are funding and repo haircuts, the gap between
the bond's price and the auction price, delivery options, and the mark-to-market if
the basis widens.
\iqlookfor{the package, funding, and the residual risks.}
\end{solution}

\begin{solution}{iq:m2:credit-default-swaps:6}
Store trades with their standard terms; nightly, bootstrap hazard curves per name and
recovery from market quotes, value each trade and aggregate risk by name and tenor, in
parallel across names; on a \omterm{def:m2:credit-default-swaps:event}{credit event}, watch the committee's decisions and the
auction result, then compute each trade's settlement at the final price, net by
counterparty and through the clearing house, and book the payments; test against
past auctions.
\iqlookfor{curves per name, parallel valuation, the event workflow.}
\end{solution}
